% ============================================================================
%  Relatorio Tecnico - Trabalho 1 de Programacao Paralela
%  Formato exigido:
%    pag. 1  -> coluna DUPLA, margens 1cm, fonte 10, somente texto
%    pag. 2  -> coluna SIMPLES: tabelas (tempo, speed-up, eficiencia) + graficos
%    pag. 3+ -> coluna SIMPLES: codigo fonte
%  Compilar com pdflatex (ou colar no Overleaf).
% ============================================================================
\documentclass[10pt,a4paper,twocolumn]{article}

\usepackage[utf8]{inputenc}
\usepackage[T1]{fontenc}
\usepackage[brazil]{babel}
\usepackage[margin=1cm]{geometry}
\usepackage{graphicx}
\usepackage{booktabs}
\usepackage{listings}
\usepackage{xcolor}
\usepackage{multicol}

\setlength{\parindent}{0pt}
\setlength{\parskip}{3pt}
\setlength{\columnsep}{6mm}
\pagestyle{empty}

% ---- Estilo para a listagem de codigo C ------------------------------------
\lstdefinestyle{codigoC}{
  language=C, basicstyle=\ttfamily\scriptsize, numbers=left,
  numberstyle=\tiny\color{gray}, keywordstyle=\color{blue!70!black},
  commentstyle=\itshape\color{green!40!black}, stringstyle=\color{red!60!black},
  showstringspaces=false, breaklines=true, tabsize=4,
  columns=fullflexible, lineskip=-1pt,
  morekeywords={\#pragma,omp,parallel,for,schedule,reduction,shared,default}
}

\begin{document}

% ============================ CABECALHO =====================================
\twocolumn[
\begin{center}
  {\large\bfseries Paralelização do Conjunto de Mandelbrot com OpenMP}\\[2pt]
  {\small Trabalho 1 --- Programação Paralela --- Grupo 20 (conta LAD \texttt{cp32020})}\\[1pt]
  {\small NOME DO INTEGRANTE 1 \quad NOME DO INTEGRANTE 2 \quad NOME DO INTEGRANTE 3}\\[1pt]
  {\small Escola Politécnica --- PUCRS \quad --- \quad \today}
\end{center}
\vspace{4mm}
]

% ======================== PAGINA 1: TEXTO ===================================

\textbf{1. O problema e seu potencial de paralelização.}
% TODO: descrever o Mandelbrot, o kernel z_{n+1}=z_n^2+c, o criterio de escape
% |z|>2 e o limite max_iter. Enfatizar:
%  - independencia total entre pixels (nenhuma dependencia de dados);
%  - eixos possiveis de paralelizacao: pixel, linha, bloco de linhas;
%  - escolha do laco de LINHAS (granularidade grossa, escrita em faixas
%    disjuntas do vetor, sem falso compartilhamento entre threads);
%  - custo por pixel MUITO variavel -> desbalanceamento intrinseco de carga,
%    que e' o ponto central deste trabalho.

\textbf{2. Ambiente de execução e metodologia.}
% TODO preencher com os dados de results/maquina.txt:
%  - nodo do cluster atlantica (LAD/PUCRS), modelo da CPU, N nucleos fisicos;
%  - gcc versao X, flags -O2 -std=c99 [-fopenmp];
%  - T(1) medido com o binario SEQUENCIAL (compilado SEM -fopenmp), 1 nucleo,
%    a partir do MESMO fonte e das MESMAS flags -O2: sem isso o speed-up sairia
%    inflado por comparar um sequencial pior compilado com um paralelo melhor;
%  - OMP_PROC_BIND=close e OMP_PLACES=cores (sem migracao, sem hyper-threads);
%  - cada ponto = MEDIANA de 3 execucoes (a mediana resiste a uma execucao
%    contaminada por outro processo no no, a media nao);
%  - job submetido com sbatch --exclusive: sem exclusividade outro grupo pode
%    dividir o no e os tempos perdem sentido;
%  - validacao por checksum: a soma das iteracoes e' identica em todas as
%    versoes e contagens de threads, garantindo corretude do paralelismo.

\textbf{3. Diretivas OpenMP utilizadas.}
% TODO: explicar cada clausula aplicada em calcula_mandelbrot():
%  - #pragma omp parallel for  no laco de LINHAS (e nao no de colunas: menos
%    aberturas de regiao paralela e granularidade maior);
%  - default(none): forca a classificacao explicita de cada variavel;
%  - reduction(+:soma): o checksum e' o unico acumulador compartilhado; a
%    reducao evita uma secao critica que serializaria o laco;
%  - variaveis do laco interno declaradas no escopo interno -> privadas por
%    construcao (mais seguro que listar em private());
%  - schedule(runtime) + omp_set_schedule() chamado a partir da linha de
%    comando (-s/-c): permite comparar static/dynamic/guided e varios chunks
%    com o BINARIO IDENTICO, sem recompilar, o que torna a Tabela 3 uma
%    comparacao justa (nenhuma diferenca de codigo gerado entre as medicoes).
% Otimizacoes: |z|^2 > 4 no lugar de sqrt; reuso de zx2/zy2 (2 multiplicacoes
% por iteracao); kernel inline e sem acesso a memoria no laco quente.

\textbf{4. Análise dos resultados --- escalabilidade forte.}
% TODO: NAO repetir numeros do grafico; EXPLICAR:
%  - por que static satura/perde eficiencia: as linhas centrais atravessam o
%    interior do conjunto e custam max_iter por pixel, enquanto as linhas das
%    bordas escapam cedo. Com blocos contiguos, as threads que recebem o meio
%    da imagem trabalham muito mais e as demais ficam ociosas no barreira
%    implicita do "for" -> o tempo total e' ditado pela thread mais lenta;
%  - por que dynamic,1 / guided recuperam a eficiencia: redistribuem linhas sob
%    demanda, ao custo de sincronizacao na fila de tarefas;
%  - de onde vem a perda residual em p alto (fracao serial da alocacao/inicio
%    da regiao paralela pela lei de Amdahl, contencao de memoria, granularidade
%    final insuficiente quando ha poucas linhas por thread).

\textbf{5. Análise --- escalabilidade fraca.}
% TODO: explicar que a carga por thread e' mantida constante escalando as duas
% dimensoes por sqrt(p) sobre a MESMA regiao do plano, o que preserva a
% distribuicao estatistica do custo por pixel. Discutir por que a eficiencia
% fraca se sustenta melhor (ou nao) que a forte: mais linhas por thread ->
% melhor granularidade e menor peso relativo da parte serial; em contrapartida
% cresce a pressao sobre memoria/cache (imagem de p x 1,92 Mpixels).

\textbf{6. Balanceamento de carga: escalonadores e chunks.}
% TODO: comentar a Tabela 3 explicando o compromisso entre desbalanceamento
% (chunk grande) e sobrecusto de escalonamento (chunk pequeno demais), e
% justificar a politica escolhida como padrao.

\textbf{7. Uso de ferramentas de IA.}
% TODO: declarar honestamente como a IA foi usada (o enunciado exige):
% ex. "o Claude Code foi usado para estruturar o codigo instrumentado, os
% scripts de medicao e a formatacao do relatorio, e para revisar o codigo do
% grupo (onde apontou um erro de compilacao com -std=c11 estrito); o grupo
% definiu o problema, validou os resultados e escreveu a analise."

\textbf{8. Conclusão.}
% TODO: 3-4 linhas.

% ==================== PAGINA 2: TABELAS E GRAFICOS ==========================
\onecolumn
\newpage

\begin{center}\textbf{\large Tabelas e gráficos}\end{center}

\textbf{Tabela 1 --- Escalabilidade forte.} Imagem fixa de
4096$\times$4096 pixels, \texttt{max\_iter} = 2000, \texttt{schedule(dynamic,1)}.
$T(1)$ é o binário \emph{sequencial} (compilado sem \texttt{-fopenmp}), em um
único núcleo. Mediana de 3 execuções. $S(p) = T_{seq}/T(p)$ e
$E(p) = S(p)/p$.\\[2pt]
\begin{tabular}{lrrrr}
\toprule
Versão & Threads & Tempo (s) & Speed-up & Eficiência (\%) \\
\midrule
sequencial & 1  & & --- & --- \\
paralela   & 1  & & & \\
paralela   & 2  & & & \\
paralela   & 4  & & & \\
paralela   & 8  & & & \\
paralela   & 16 & & & \\
% acrescentar linhas ate o total de nucleos do no
\bottomrule
\end{tabular}

\vspace{6mm}
\textbf{Tabela 2 --- Escalabilidade fraca.} Carga por thread constante:
1024$\times$1024 pixels por thread, lado escalado por $\sqrt{p}$ sobre a
\emph{mesma} janela do plano complexo. Aqui
$S(p) = p \cdot T_{seq}/T(p)$ e $E(p) = T_{seq}/T(p)$ --- ou seja, a
eficiência mede o quanto o tempo conseguiu permanecer constante.\\[2pt]
\begin{tabular}{rlrrr}
\toprule
Threads & Imagem & Mpixels & Tempo (s) & Eficiência (\%) \\
\midrule
1  & 1024$\times$1024 & 1,05  & & \\
2  & 1448$\times$1448 & 2,10  & & \\
4  & 2048$\times$2048 & 4,19  & & \\
8  & 2896$\times$2896 & 8,39  & & \\
16 & 4096$\times$4096 & 16,78 & & \\
\bottomrule
\end{tabular}

\vspace{6mm}
\textbf{Tabela 3 --- Balanceamento de carga: escalonador e granularidade.}
Todos os núcleos do nó, imagem 4096$\times$4096.\\[2pt]
\begin{tabular}{llrr}
\toprule
\texttt{schedule} & Chunk & Tempo (s) & Relativo ao melhor \\
\midrule
\texttt{static}  & padrão & & \\
\texttt{static}  & 1  & & \\
\texttt{static}  & 8  & & \\
\texttt{static}  & 64 & & \\
\texttt{dynamic} & 1  & & \\
\texttt{dynamic} & 8  & & \\
\texttt{dynamic} & 64 & & \\
\texttt{guided}  & padrão & & \\
\texttt{guided}  & 1  & & \\
\bottomrule
\end{tabular}

\vspace{8mm}
% Graficos gerados na planilha fornecida pelo professor, exportados como PNG/PDF.
% \includegraphics[width=0.9\linewidth]{grafico_forte.png}
% \includegraphics[width=0.9\linewidth]{grafico_fraca.png}

% ==================== PAGINA 3+: CODIGO FONTE ===============================
\newpage
\begin{center}\textbf{\large Código fonte --- \texttt{mandelbrot.c}}\end{center}
\lstinputlisting[style=codigoC]{../mandelbrot.c}

\newpage
\begin{center}\textbf{\large Script de medição --- \texttt{medir.sh}}\end{center}
\lstinputlisting[style=codigoC,language=bash]{../medir.sh}

\newpage
\begin{center}\textbf{\large Análise dos resultados --- \texttt{analisar.sh}}\end{center}
\lstinputlisting[style=codigoC,language=bash]{../analisar.sh}

\newpage
\begin{center}\textbf{\large Submissão ao Slurm --- \texttt{job.sh}}\end{center}
\lstinputlisting[style=codigoC,language=bash]{../job.sh}

\end{document}
